\documentclass[11pt,a4paper]{article}

% ---------------------------------------------------------
% PREAMBLE (stable, warning-free, Overleaf-safe)
% ---------------------------------------------------------
\usepackage[utf8]{inputenc}
\usepackage[T1]{fontenc}
\usepackage{lmodern}
\usepackage{amsmath, amssymb, amsfonts}
\usepackage{bm}
\usepackage{geometry}
\usepackage{graphicx}
\usepackage{hyperref}
\usepackage{cite}
\usepackage{authblk}
\usepackage{physics}
\usepackage{mathtools}

\geometry{margin=2.4cm}

\title{\textbf{Companion LXXXII: Transport-Based Parameter Inference and $\alpha$-Sensitivity Maps for the Relaxation-Driven Cyclic Cosmology}}
\author{Michael Lehmann \\ \texttt{mi.lehmann@gmx.de}}
\date{April 2026}

\begin{document}
\maketitle

\begin{abstract}
Transport-based statistics provide a natural way to compare distributions of 
persistence diagrams across cosmological models. Building on the diffusion--optimal 
transport framework developed in previous Companions, this work introduces 
transport-based parameter inference and $\alpha$-sensitivity maps for the 
Relaxation-Driven Cyclic Cosmology (RDCC). The relaxation parameter $\alpha$ 
controls the scale-dependent suppression of the gravitational potential, which in 
turn deforms the distribution of weak-lensing persistence diagrams. We construct 
transport distances and neural transport maps between RDCC and $\Lambda$CDM 
ensembles, and define sensitivity fields that quantify how these transport 
statistics respond to variations in $\alpha$. The resulting framework provides a 
direct link between geometric transport quantities and observable constraints on 
$\alpha$ from Euclid-like weak-lensing surveys.
\end{abstract}
% ---------------------------------------------------------
\section{Introduction}

Weak-lensing convergence maps encode the projected large-scale structure of the 
Universe. Their topological structure can be quantified using persistence 
diagrams, which capture the birth and death of connected components, loops, and 
voids. Previous Companions developed a hierarchy of geometric tools for these 
diagrams, including Wasserstein distances, entropic optimal transport, 
Schrödinger bridges, diffusion models, and neural transport flows.

In the Relaxation-Driven Cyclic Cosmology (RDCC), the relaxation parameter 
$\alpha$ controls the scale-dependent suppression of the gravitational potential. 
This induces characteristic deformations in the distribution of persistence 
diagrams relative to $\Lambda$CDM. The goal of this Companion is to translate 
these deformations into \emph{transport-based parameter constraints} and explicit 
\emph{$\alpha$-sensitivity maps} that can be evaluated on Euclid-like mock 
catalogues.
% ---------------------------------------------------------
\section{Transport Distances as Summary Statistics}

Let $\mathcal{D}_{\Lambda}$ and $\mathcal{D}_{\alpha}$ denote ensembles of 
persistence diagrams from $\Lambda$CDM and RDCC with parameter $\alpha$, 
respectively. We define a transport-based discrepancy

\begin{equation}
    \mathcal{T}(\alpha)
    =
    \mathbb{E}_{D \sim \mathcal{D}_{\Lambda}}
    \left[
        W_{p}\bigl(D, T_{\alpha}(D)\bigr)
    \right],
\end{equation}

where $W_{p}$ is a Wasserstein or Sinkhorn distance, and $T_{\alpha}$ is a neural 
transport map pushing $\mathcal{D}_{\Lambda}$ towards $\mathcal{D}_{\alpha}$.

The function $\mathcal{T}(\alpha)$ acts as a \emph{summary statistic}: it 
vanishes when the two ensembles are indistinguishable and grows as RDCC deviates 
from $\Lambda$CDM. Its scale dependence can be probed by restricting diagrams to 
features in specific birth--death or angular multipole ranges.
% ---------------------------------------------------------
\section{\texorpdfstring{$\alpha$}{alpha}-Sensitivity Maps}

We define the $\alpha$-sensitivity of the transport statistic as

\begin{equation}
    S_{\alpha}(\ell)
    =
    \frac{\partial}{\partial \alpha}
    \mathcal{T}_{\ell}(\alpha)
    \bigg|_{\alpha = \alpha_{0}},
\end{equation}

where $\mathcal{T}_{\ell}(\alpha)$ is the transport discrepancy computed from 
features associated with angular scale $\ell$, and $\alpha_{0}$ is a fiducial 
value.

In practice, $S_{\alpha}(\ell)$ is estimated by finite differences using RDCC 
mocks at nearby $\alpha$ values. The resulting function provides an 
\emph{$\ell$-resolved sensitivity map} indicating which scales contribute most 
strongly to constraining $\alpha$.
% ---------------------------------------------------------
\section{Likelihood-Free Inference with Transport Statistics}

Given observed persistence diagrams $\mathcal{D}_{\mathrm{obs}}$ from a 
Euclid-like survey, we define an empirical transport discrepancy

\begin{equation}
    \widehat{\mathcal{T}}(\alpha)
    =
    \mathbb{E}_{D \sim \mathcal{D}_{\mathrm{obs}}}
    \left[
        W_{p}\bigl(D, T_{\alpha}^{\mathrm{mock}}(D)\bigr)
    \right],
\end{equation}

where $T_{\alpha}^{\mathrm{mock}}$ is trained on simulations at parameter 
$\alpha$. This quantity can be used in a likelihood-free inference scheme:

\begin{itemize}
    \item sample $\alpha$ from a prior,
    \item generate RDCC mocks and train $T_{\alpha}^{\mathrm{mock}}$,
    \item compute $\widehat{\mathcal{T}}(\alpha)$,
    \item accept or weight $\alpha$ according to the discrepancy.
\end{itemize}

The resulting posterior on $\alpha$ is driven directly by transport-based 
topological differences between RDCC and $\Lambda$CDM.
% ---------------------------------------------------------
\section{Connection to Previous Geometric Companions}

The present framework leverages several structures developed earlier:

\begin{itemize}
    \item Wasserstein and Sinkhorn distances (Companions LXXVI--LXXVII),
    \item Schrödinger bridges and diffusion-regularized transport (LXXVIII),
    \item score-based diffusion models (LXXIX),
    \item neural Schrödinger flows and unified diffusion--OT geometry (LXXX--LXXXI).
\end{itemize}

Neural transport maps $T_{\alpha}$ can be instantiated as:

\begin{itemize}
    \item the mean flow of a neural Schrödinger SDE,
    \item the deterministic component of a unified diffusion--OT model,
    \item or an explicit OT map parameterized by a neural network.
\end{itemize}

In all cases, the RDCC parameter $\alpha$ enters through the modified convergence 
field and the resulting deformation of the persistence-diagram distribution.
% ---------------------------------------------------------
\section{Conclusion}

Transport-based statistics provide a natural and geometrically meaningful way to 
compare persistence-diagram ensembles across cosmological models. In the 
Relaxation-Driven Cyclic Cosmology, the relaxation parameter $\alpha$ induces 
scale-dependent deformations that can be captured by neural transport maps and 
Wasserstein or Sinkhorn discrepancies. The $\alpha$-sensitivity maps introduced 
here quantify which angular scales and topological features contribute most 
strongly to constraining $\alpha$ in Euclid-like weak-lensing surveys. The present 
Companion thus connects the geometric machinery of previous Companions directly 
to observable parameter inference, addressing the need for explicit links between 
theory and constraints.

\bibliographystyle{apsrev4-2}
\nocite{*}
\bibliography{references}


\end{document}
